\begin{figure}[ht]
\centering
\begin{tikzpicture}[x=1cm, y=1cm, font=\footnotesize,
  q/.style={draw=engblue, thick, fill=engblue!8, rounded corners=2pt,
            minimum width=4cm, minimum height=0.9cm, align=center},
  d/.style={draw=engred, thick, fill=engred!10, rounded corners=2pt,
            minimum width=3.4cm, minimum height=0.9cm, align=center},
  arr/.style={->, thick}
]

\node[q] (state) at (0,    7) {1. State the measurand\\(quantity, units, $u_c$, range, rate)};
\node[q] (princ) at (0,  5.4) {2. Identify candidate sensing principles\\
                                from \S 5.2 catalogue};
\node[q] (surv)  at (0,  3.8) {3. Survey 2--4 commercial sensors per principle\\
                                read datasheets in full};
\node[q] (cmp)   at (0,  2.2) {4. Compare specifications against requirements\\
                                identify marginal specs};
\node[q] (env)   at (-3, 0.4) {5. Environmental constraints\\
                                (T, RH, vibr., EMI, chemistry)};
\node[q] (intg)  at (3,  0.4) {6. Integration constraints\\
                                (signal, ADC, mounting, cable)};
\node[q] (cost)  at (0, -1.4) {7. Cost and supply\\
                                (price, lead time, stock)};
\node[d] (sel)   at (0, -3.0) {8. Document the selection\\
                                (record survives the engineer)};

\draw[arr] (state) -- (princ);
\draw[arr] (princ) -- (surv);
\draw[arr] (surv)  -- (cmp);
\draw[arr] (cmp)   -- (env);
\draw[arr] (cmp)   -- (intg);
\draw[arr] (env)   -- (cost);
\draw[arr] (intg)  -- (cost);
\draw[arr] (cost)  -- (sel);

% loopback: failed candidate returns to step 2
\draw[arr, dashed, color=engred] (env.west) -- ++(-1.6, 0)
  |- (princ.west) node[midway, font=\scriptsize\itshape, color=engred,
     anchor=south, rotate=90, xshift=15pt] {if eliminated};

\end{tikzpicture}
\caption[The sensor selection procedure]{Sensor selection as a
checklist that becomes judgment with use. Step 1 forces the
measurand into the same place as the budget. Steps 2 to 4 narrow
the field by principle and by specification. Steps 5 and 6 enforce
the environment and the integration, which eliminate
perfect-on-the-bench candidates that cannot survive the field.
Step 7 names the trade-offs the engineer accepts (lead time, stock,
price). Step 8 commits the selection to the design record, where
it survives the engineer's tenure and informs later substitutions.
The dashed return from steps 5 and 6 to step 2 is the working
case: a candidate eliminated by environment or integration sends
the engineer back to the principle catalogue.}
\label{fig:vol01:ch05:selection-tree}
\end{figure}
